\documentclass{beamer}
\usetheme{Madrid}
\usecolortheme{whale}

\title[CVD Risk Classification]{Cardiovascular Disease Risk Classification Pipeline}
\subtitle{Capstone Project -- Semester S4}
\author{Assadiki Zaid}
\institute{École Hassania des Travaux Publics (EHTP)}
\date{June 18, 2026}

\begin{document}

\frame{\titlepage}

\begin{frame}{Slide 1: Problem Motivation \& Framing}
    \begin{itemize}
        \item \textbf{Clinical Challenge:} Cardiovascular disease is the leading cause of global mortality.
        \item \textbf{Objective:} Build an end-to-end data pipeline to classify patients as high-risk ($cardio=1$) or low-risk ($cardio=0$).
        \item \textbf{Data Ingestion:} 70,000 observations containing 11 clinical features (biometrics, lab tests, lifestyle habits) from Kaggle.
    \end{itemize}
\end{frame}

\begin{frame}{Slide 2: Pipeline Design}
    \begin{block}{Architecture Highlights}
        \begin{enumerate}
            \item Ingestion $\rightarrow$ Ingest script descriptions.
            \item Wrangling $\rightarrow$ Outlier comparison (IQR vs Z-Score) and physical thresholds.
            \item Feature Engineering $\rightarrow$ Derive BMI, Pulse Pressure, and Age-Cholesterol Risk.
            \item SQLite Storage $\rightarrow$ Store cleaned data for analytical verification.
            \item Preprocessing $\rightarrow$ ColumnTransformers to enforce ``No Data Leakage''.
            \item Modeling $\rightarrow$ Fit pipelines for LR, KNN, Random Forest, and XGBoost.
        \end{enumerate}
    \end{block}
\end{frame}

\begin{frame}{Slide 3: Exploratory Highlights \& Surprising Insight}
    \begin{alertblock}{The Pulse Pressure Discovery}
        While we expected raw Systolic BP (\textit{ap\_hi}) to be the primary classifier, EDA revealed that the engineered feature \textbf{Pulse Pressure} ($\textit{ap\_hi} - \textit{ap\_lo}$) provides a significantly cleaner risk boundary, particularly in younger cohorts (under 55 years old).
    \end{alertblock}
    \begin{itemize}
        \item Pulse Pressure acts as a proxy for early-stage arterial stiffness.
        \item Reflects dynamic cardiac load more accurately than static blood pressure readings.
    \end{itemize}
\end{frame}

\begin{frame}{Slide 4: Model Benchmarking}
    \begin{table}
        \begin{tabular}{llllll}
            \hline
            Model & Accuracy & Precision & Recall & $F_1$-score & ROC-AUC \\
            \hline
            \textbf{Tuned XGBoost} & \textbf{73.57\%} & \textbf{75.39\%} & \textbf{70.21\%} & \textbf{72.71\%} & \textbf{80.12\%} \\
            Logistic Regression & 72.84\% & 74.92\% & 68.61\% & 71.63\% & 79.28\% \\
            Random Forest & 71.21\% & 72.10\% & 69.15\% & 70.59\% & 77.40\% \\
            K-Nearest Neighbors & 69.80\% & 70.35\% & 68.10\% & 69.21\% & 74.50\% \\
            \hline
        \end{tabular}
    \end{table}
    \begin{itemize}
        \item XGBoost optimized via GridSearchCV (\texttt{learning\_rate=0.05}, \texttt{max\_depth=3}).
        \item Prioritized $F_1$-score over accuracy to protect patients (reducing fatal false negatives) and save clinic resources (reducing false positives).
    \end{itemize}
\end{frame}

\begin{frame}{Slide 5: Clinical Recommendations}
    \begin{itemize}
        \item \textbf{Deployment:} Integrate the serialized model (\texttt{best\_cardio\_model.joblib}) into the EHR system as a triage decision-support tool.
        \item \textbf{Operational Safeguards:} Human-in-the-loop validation; restrict model usage to ages 30--65.
        \item \textbf{Clinical Impact:} Proactive therapeutic intervention, reduced hospital readmission rates, and optimal patient scheduling.
    \end{itemize}
\end{frame}

\end{document}
